\section{Cinemática Inversa}
\label{sec_cinematica_inversa}

Dados a pose desejada do efetuador (posição e orientação), o problema de cinemática inversa consiste em determinar todos os possíveis conjuntos de ângulos de junta que compõem a pose.
Na cinemática de manipuladores, essa determinação da pose do efetuador é um problema fundamental.
O problema é não linear, pois envolve equações trigonométricas resultantes da composição de transformações homogêneas \cite{craig2005,spong2006robot}.
Além disso, manipuladores revolutos podem admitir múltiplas soluções articulares para uma mesma pose, exigindo que a escolha final respeite continuidade da trajetória e limites mecânicos \cite{craig2005}.

\subsection{Erro operacional}

A formulação diferencial da cinemática inversa baseia-se na definição do erro no espaço operacional.
Seja ${}^{0}\!\vec{x}_{des}$ a posição desejada do \gls{eef} e ${}^{0}\!\vec{x}$ a posição atual obtida pela cinemática direta, ambas expressas no referencial $\{0\}$.
Define-se o erro cartesiano como \cite{sciavicco_siciliano2010}:

\begin{equation}
\label{eq_erro_operacional}
\vec{e}
=
{}^{0}\!\vec{x}_{des}
-
{}^{0}\!\vec{x}.
\end{equation}

Esse erro representa a diferença instantânea entre a posição desejada e a posição atual do efetuador e constitui a variável regulada pelo algoritmo de cinemática inversa diferencial.


\subsection{Erro de orientação}

Além da posição, é necessário considerar o alinhamento do efetuador com uma direção desejada no espaço, particularmente na fase de aproximação para acoplamento.
Seja ${}^{0}\!\hat{z}_{E}$ o versor associado ao eixo $\hat{z}$ do efetuador $\{E\}$, expresso no referencial $\{0\}$, e seja ${}^{0}\!\hat{d}$ um versor que representa a direção desejada, também expresso em $\{0\}$.
Define-se o erro de orientação como:

\begin{equation}
\label{eq_erro_orientacao}
\vec{e}_{\omega}
=
{}^{0}\!\hat{z}_{E}
\times
{}^{0}\!\hat{d}.
\end{equation}

De modo que $\vec{e}_{\omega}=\vec{0}$ indica alinhamento entre o eixo do efetuador e a direção desejada.
Esse erro vetorial é utilizado na construção da componente angular da velocidade operacional desejada.

%%%%%%%%%%%%%%%%%%%%%%%%%%%%%%%%%%%%%%%%%%%%%%%%%%%%%%%%%%

\subsection{Velocidade cartesiana desejada}

Seguindo a formulação de controle cinemático apresentada em \cite{sciavicco_siciliano2010}, utiliza-se uma lei proporcional no espaço operacional, separando as componentes linear e de velocidade angular:

\begin{equation}
\vec{v}_{des}
=
\dot{\vec{x}}_{des}
+
K_{p}\,\vec{e},
\label{eq_vd_desejada}
\end{equation}

\begin{equation}
\vec{\omega}_{des}
=
K_{\omega}\,\vec{e}_{\omega}.
\label{eq_omegad_desejada}
\end{equation}

Em que $\dot{\vec{x}}_{des}$ é a velocidade cartesiana de referência.
Os ganhos $K_{p}>0$ e $K_{\omega}>0$ são escalares e ponderam, respectivamente, os termos proporcionais associados aos erros de posição e de orientação.
A velocidade operacional desejada é reunida no vetor:

\begin{equation}
\label{eq_velocidade_operacional_desejada}
\vec{y}_{des}
=
\begin{bmatrix}
\vec{v}_{des} \\
\vec{\omega}_{des}
\end{bmatrix}.
\end{equation}

O qual combina o objetivo de posicionamento com o objetivo de alinhamento do efetuador.
Essa velocidade operacional desejada é então mapeada para o espaço das juntas por meio da Jacobiana do manipulador, resultando nas velocidades articulares desejadas utilizadas na cinemática inversa diferencial.


%%%%%%%%%%%%%%%%%%%%%%%%%%%%

\subsection{Lei de cinemática diferencial inversa}

A relação entre a velocidade operacional desejada e as velocidades articulares é obtida por meio da pseudo-inversa amortecida da Jacobiana geométrica do manipulador.
Seja $\vec{\xi}$ o vetor de coordenadas articulares e $\dot{\vec{\xi}}$ o vetor de velocidades articulares.
A Jacobiana geométrica $\mathbf{J}_{g}(\vec{\xi})$ relaciona $\dot{\vec{\xi}}$ à velocidade operacional do efetuador, sendo:

\begin{equation}
\mathbf{J}_{g}(\vec{\xi})
=
\begin{bmatrix}
\mathbf{J}_{v}(\vec{\xi}) \\
\mathbf{J}_{\omega}(\vec{\xi})
\end{bmatrix}.
\label{eq_jacobiana_particionada}
\end{equation}

A lei de cinemática inversa diferencial é escrita como

\begin{equation}
\dot{\vec{\xi}}
=
\mathbf{J}_{\lambda}^{\dagger}(\vec{\xi})\,\vec{y}_{des}.
\label{eq_ik_diferencial}
\end{equation}

Em que $\vec{y}_{des}$ é a velocidade operacional desejada definida em \eqref{eq_velocidade_operacional_desejada} e
$\mathbf{J}_{\lambda}^{\dagger}(\vec{\xi})$ denota a pseudo-inversa amortecida, com parâmetro de amortecimento $\lambda>0$, definida em \eqref{eq_pseudoinversa_amortecida2}.



%%%%%%%%%%%%%%%%%%%%%%%%%%%%

\subsection{Pseudo-inversa amortecida}

A pseudo-inversa amortecida empregada na lei de cinemática diferencial inversa \eqref{eq_ik_diferencial} é definida segundo \cite{sciavicco_siciliano2010}, como:

\begin{equation}
\label{eq_pseudoinversa_amortecida2}
\mathbf{J}_{\lambda}^{\dagger}
=
\mathbf{J}^{T}
\bigl(
\mathbf{J}\,\mathbf{J}^{T}
+
\lambda^{2}\mathbf{I}_{m}
\bigr)^{-1}.
\end{equation}

O termo $\lambda>0$ é o fator de amortecimento empregado para regularizar a inversão de $\mathbf{J}\mathbf{J}^{T}$, $\mathbf{I}_{m}$ é a matriz identidade de ordem $m$ (número de coordenadas da tarefa) e $\mathbf{J}$ representa a Jacobiana da tarefa considerada (nesta seção, $\mathbf{J}=\mathbf{J}_{g}(\vec{\xi})$).
Essa pseudo-inversa amortecida melhora a robustez numérica e o comportamento na vizinhança de singularidades.

%%%%%%%%%%%%%%%%%%%%%%%%%%%%%%%

\subsection{Integração da lei diferencial}

Uma vez determinada a velocidade articular instantânea $\dot{\vec{\xi}}(t)$ pela lei de cinemática diferencial inversa, a trajetória no espaço de juntas é obtida por integração explícita no tempo.
Utilizando o método de Euler explícito com passo $\Delta t$, tem-se:

\begin{equation}
\label{eq_integracao_explicita_xi}
\vec{\xi}(t + \Delta t)
=
\vec{\xi}(t)
+
\dot{\vec{\xi}}(t)\,\Delta t.
\end{equation}

Essa atualização é aplicada iterativamente, de modo que, a cada passo de integração, o erro operacional é recalculado e a lei diferencial de cinemática inversa é novamente avaliada, conforme proposto em \cite{sciavicco_siciliano2010}.
